\section{Conception détaillée}

\subsection{Identité, demandes de démonstration et invitations}

Trois parcours créent des comptes, chacun avec une règle différente :

\begin{enumerate}[leftmargin=8mm]
  \item un visiteur peut créer directement un compte \textbf{client}. L'adresse doit
  être vérifiée avant la connexion locale ;
  \item un client peut utiliser Google OAuth. L'identité Google est liée à un compte
  local, mais ne fournit ni rôle ni organisation ;
  \item un administrateur, opérateur ou technicien est toujours créé par invitation.
  Le destinataire active lui-même son compte et choisit son mot de passe.
\end{enumerate}

La demande de démonstration constitue l'entrée commerciale d'une organisation. Le
visiteur renseigne ses coordonnées, la société, la taille du réseau et plusieurs
objectifs opérationnels. Une confirmation est envoyée au demandeur sans lien
interne. Le super administrateur ouvre ensuite la demande, la rejette ou lance sa
revue. Une acceptation provisionne, dans une transaction, l'organisation, son
administrateur en état \texttt{pending}, l'essai commercial et une invitation
temporaire. Le compte ne devient actif qu'après acceptation du lien à usage unique.

\begin{figure}[H]
  \centering
  \resizebox{0.96\textwidth}{!}{\begin{tikzpicture}[
  font=\sffamily\scriptsize,
  participant/.style={draw=ctline, rounded corners=1.5mm, fill=white,
    minimum width=27mm, minimum height=9mm, align=center},
  call/.style={-{Latex[length=2mm]}, thick, draw=ctgray},
  return/.style={-{Latex[length=2mm]}, dashed, draw=ctgreen}
]
  \node[participant] (visitor) at (0,0) {Visiteur};
  \node[participant] (react) at (3.2,0) {React};
  \node[participant] (laravel) at (6.4,0) {Laravel};
  \node[participant] (db) at (9.6,0) {PostgreSQL};
  \node[participant] (mail) at (12.8,0) {Service email};
  \node[participant] (admin) at (16,0) {Super admin};

  \foreach \x in {0,3.2,6.4,9.6,12.8,16}
    \draw[densely dashed, ctline] (\x,-0.5) -- (\x,-6.2);

  \draw[call] (0,-0.9) -- node[above]{Envoyer la demande} (3.2,-0.9);
  \draw[call] (3.2,-1.25) -- node[above]{POST \texttt{/demo-requests}} (6.4,-1.25);
  \draw[call] (6.4,-1.6) -- node[above]{Valider et enregistrer} (9.6,-1.6);
  \draw[call] (6.4,-2.0) -- node[above]{Accusé de réception} (12.8,-2.0);
  \draw[return] (12.8,-2.35) -- (0,-2.35);
  \draw[call] (16,-2.8) -- node[above]{Lancer la revue} (6.4,-2.8);
  \draw[call] (16,-3.2) -- node[above]{Accepter et provisionner} (6.4,-3.2);
  \draw[call] (6.4,-3.65) -- node[above, align=center]{Transaction : organisation, essai,\\admin pending, invitation hashée} (9.6,-3.65);
  \draw[call] (6.4,-4.25) -- node[above]{Envoyer le lien d'activation} (12.8,-4.25);
  \draw[return] (12.8,-4.65) -- node[above]{Lien temporaire} (0,-4.65);
  \draw[call] (0,-5.1) -- node[above]{Choisir le mot de passe} (6.4,-5.1);
  \draw[call] (6.4,-5.55) -- node[above]{Activer et consommer le token} (9.6,-5.55);
  \draw[return] (6.4,-5.95) -- node[above]{Accès à l'onboarding} (3.2,-5.95);
\end{tikzpicture}
}
  \caption{Provisionnement d'une organisation à partir d'une demande de démonstration}
  \label{fig:account-provisioning-sequence}
\end{figure}

Les invitations suivent les invariants suivants :

\begin{itemize}[leftmargin=7mm]
  \item le token aléatoire n'est envoyé qu'une fois et seul son condensat SHA-256 est
  stocké ;
  \item le lien possède une expiration et devient inutilisable après acceptation,
  révocation ou remplacement ;
  \item un rappel renouvelle le token, car sa valeur originale ne peut pas être
  récupérée ;
  \item un administrateur ne peut inviter que des opérateurs ou techniciens de son
  organisation ;
  \item l'utilisateur invité choisit son secret ; l'administrateur ne connaît jamais
  son mot de passe.
\end{itemize}

\subsection{Mise en service d'une station}

Le commissioning rassemble la station, son emplacement, son matériel, ses
connecteurs et sa cible OCPP dans un assistant unique. La création est atomique :
si une référence, une identité OCPP ou un connecteur est invalide, aucune partie
incomplète n'est conservée.

\begin{table}[H]
  \centering
  \caption{Cibles de mise en service OCPP}
  \label{tab:commissioning-targets}
  \begin{tabularx}{\textwidth}{@{}p{3.5cm}Y Y@{}}
    \toprule
    \textbf{Cible} & \textbf{Comportement} & \textbf{Usage}\\
    \midrule
    Borne physique ou externe &
    Génération d'un secret Basic Auth individuel de 48 caractères, affiché une
    seule fois et stocké uniquement sous forme de hash. &
    Intégration future d'un chargeur réel ou d'un simulateur tiers.\\
    Simulateur SAP local &
    Création de l'inventaire puis ajout du profil au manifeste local par une commande
    Artisan dédiée. &
    Tests reproductibles du stage.\\
    Inventaire seul &
    Création sans identité ni secret OCPP. L'état reste non disponible jusqu'au
    provisioning. &
    Préparation administrative d'un futur actif.\\
    \bottomrule
  \end{tabularx}
\end{table}

La référence interne et l'identité OCPP peuvent différer, mais leur égalité simplifie
l'exploitation. Le libellé \texttt{A1} d'un connecteur est destiné à l'utilisateur ;
\texttt{ocpp\_connector\_id} est l'entier réellement transmis par la borne. Pour le
simulateur SAP, ces identifiants doivent être contigus à partir de 1. Après
provisionnement, la borne se connecte à
\texttt{/ocpp/<identité>} avec le sous-protocole \texttt{ocpp1.6}.

\subsection{Ingestion OCPP et disponibilité calculée}

\subsubsection{Séparation entre preuve brute et projection métier}

Un statut reçu de la borne n'est pas directement présenté comme vérité métier.
L'événement complet est d'abord enregistré dans \texttt{ocpp\_events}. Les champs
\texttt{ocpp\_*} de la station et du connecteur forment ensuite une projection
technique, tandis que les champs \texttt{status}, \texttt{availability\_reason} et
\texttt{availability\_source} constituent la projection métier.

Cette séparation apporte trois garanties :

\begin{itemize}[leftmargin=7mm]
  \item un événement tardif ne remplace pas une preuve plus récente ;
  \item les règles peuvent évoluer sans supprimer l'historique technique ;
  \item chaque changement visible peut être expliqué par une transition et son
  événement source.
\end{itemize}

\begin{figure}[H]
  \centering
  \resizebox{0.98\textwidth}{!}{\begin{tikzpicture}[
  font=\sffamily\small,
  pipelinebox/.style={draw=ctline, rounded corners=2mm, thick, fill=white,
    minimum width=29mm, minimum height=13mm, align=center},
  pipelinefocus/.style={pipelinebox, draw=ctgreen, fill=ctmint},
  flow/.style={-{Latex[length=2.2mm]}, thick, draw=ctgray}
]
  \node[pipelinebox] (station) {Borne ou\\simulateur};
  \node[pipelinebox, right=7mm of station] (gateway) {Passerelle\\OCPP};
  \node[pipelinebox, right=7mm of gateway] (event) {Événement brut\\immuable};
  \node[pipelinefocus, right=7mm of event] (project) {Projection\\technique};
  \node[pipelinefocus, right=7mm of project] (rules) {Moteur de\\disponibilité};
  \node[pipelinebox, right=7mm of rules] (business) {État métier +\\transition};
  \node[pipelinebox, below=10mm of business] (alerts) {Alerte\\dédupliquée};
  \node[pipelinebox, below=10mm of rules] (reverb) {Événement Reverb\\mis en queue};
  \node[pipelinebox, below=10mm of project] (ui) {Carte, détail,\\tableaux de bord};

  \draw[flow] (station) -- node[above, font=\scriptsize]{WS} (gateway);
  \draw[flow] (gateway) -- node[above, font=\scriptsize]{HMAC} (event);
  \draw[flow] (event) -- (project);
  \draw[flow] (project) -- (rules);
  \draw[flow] (rules) -- (business);
  \draw[flow] (business) -- (alerts);
  \draw[flow] (business) -- (reverb);
  \draw[flow] (reverb) -- (ui);
  \draw[flow] (alerts.south) -- ++(0,-5mm) -| (ui.south);
\end{tikzpicture}
}
  \caption{Chaîne de traitement d'un événement OCPP}
  \label{fig:availability-pipeline}
\end{figure}

\subsubsection{Règles de disponibilité}

La résolution applique les règles par priorité. Une décision manuelle ou une
maintenance planifiée domine toujours la télémétrie. Ensuite, une fermeture
WebSocket explicite ou 90 secondes sans message produit l'état
\texttt{offline}. Le délai représente trois périodes de heartbeat de 30 secondes.
Enfin, les états des connecteurs déterminent la capacité réelle de la station.

\begin{longtable}{@{}p{1.0cm}p{5.2cm}p{3.5cm}p{4.4cm}@{}}
  \caption{Matrice de décision de disponibilité}
  \label{tab:availability-decision-matrix}\\
  \toprule
  \textbf{Prio.} & \textbf{Preuve ou règle} & \textbf{État station} &
  \textbf{Conséquence}\\
  \midrule
  \endfirsthead
  \toprule
  \textbf{Prio.} & \textbf{Preuve ou règle} & \textbf{État station} &
  \textbf{Conséquence}\\
  \midrule
  \endhead
  1 & Désactivation manuelle & \texttt{unavailable} &
  Recharge interdite jusqu'à levée de l'override.\\
  2 & Maintenance active & \texttt{maintenance} &
  Indisponibilité planifiée et commande \texttt{ChangeAvailability}.\\
  3 & Déconnexion ou silence supérieur à 90 secondes & \texttt{offline} &
  Connecteurs hors ligne et alerte de communication.\\
  4 & Au moins un connecteur \texttt{Available} & \texttt{available} &
  Démarrage distant autorisé si l'organisation est active.\\
  5 & Aucun libre, au moins un en transaction & \texttt{charging} &
  Station vivante mais sans place libre correspondante.\\
  6 & Aucun libre, au moins un réservé & \texttt{reserved} &
  Réservation visible ; recharge publique non immédiate.\\
  7 & Aucun utilisable, au moins un \texttt{Faulted} & \texttt{faulted} &
  Alerte de défaut dédupliquée.\\
  8 & Connectée sans preuve exploitable & \texttt{unavailable} &
  Attente d'un \texttt{StatusNotification} valide.\\
  \bottomrule
\end{longtable}

Chaque variation de l'état, de la raison ou de la source crée une
\texttt{AvailabilityTransition}. Les alertes automatiques utilisent une clé de
déduplication stable : la même panne ne crée pas une série d'alertes identiques,
mais une récurrence après résolution rouvre le suivi. L'événement
\texttt{StationAvailabilityChanged} invalide les vues React concernées via Reverb.

\subsection{Parcours de recharge et transaction OCPP}

\subsubsection{Assistant côté client}

Le parcours client suit cinq étapes afin de correspondre aux actions physiques :

\begin{enumerate}[leftmargin=8mm]
  \item choisir une station publique disponible ou ouvrir directement une station ;
  \item sélectionner un connecteur compatible et libre ;
  \item brancher le câble ; l'interface attend réellement le statut OCPP
  \texttt{Preparing} et n'utilise pas une confirmation manuelle ;
  \item choisir le moyen de paiement et, facultativement, une limite d'énergie, de
  montant ou de durée ;
  \item autoriser le paiement puis demander le démarrage distant.
\end{enumerate}

Le rayon \emph{Near me} est une préférence de compte. La géolocalisation n'est
demandée qu'après une action explicite et reste dans le navigateur. Un QR de
connecteur contient un lien profond vers la station et le connecteur ; il ne
contient ni idTag, ni secret OCPP, ni jeton de paiement.

\subsubsection{Tentative, session et transaction}

Trois objets distincts évitent de facturer une énergie qui n'a pas été délivrée :

\begin{description}[leftmargin=34mm,style=nextline]
  \item[\texttt{ChargingAttempt}] intention du client, préautorisation et commande de
  démarrage ;
  \item[\texttt{OcppTransaction}] preuve technique issue de la borne ;
  \item[\texttt{ChargingSession}] session métier créée seulement après un
  \texttt{StartTransaction} accepté.
\end{description}

\begin{figure}[p]
  \centering
  \resizebox{0.96\textwidth}{!}{\begin{tikzpicture}[
  x=1cm,
  y=1cm,
  font=\sffamily\tiny,
  participant/.style={draw=ctline, rounded corners=1.5mm, fill=white,
    minimum width=19mm, minimum height=8mm, align=center},
  call/.style={-{Latex[length=1.8mm]}, thick, draw=ctgray},
  return/.style={-{Latex[length=1.8mm]}, dashed, draw=ctgreen},
  msg/.style={fill=white, inner sep=0.7pt, font=\sffamily\tiny}
]
  \node[participant] (client) at (0,0) {Client};
  \node[participant] (app) at (2.45,0) {Application};
  \node[participant] (payment) at (4.9,0) {Paiement};
  \node[participant] (gateway) at (7.35,0) {OCPP};
  \node[participant] (station) at (9.8,0) {Borne};
  \node[participant] (worker) at (12.25,0) {Worker};

  \foreach \x in {0,2.45,4.9,7.35,9.8,12.25}
    \draw[densely dashed, ctline] (\x,-0.45) -- (\x,-9.0);

  \draw[call] (0,-0.75) -- node[msg, above]{Choix station / connecteur} (2.45,-0.75);
  \draw[call] (2.45,-1.15) -- node[msg, above]{Préautoriser} (4.9,-1.15);
  \draw[return] (4.9,-1.5) -- node[msg, above]{Autorisé} (2.45,-1.5);
  \draw[call] (2.45,-1.9) -- node[msg, above]{Commande pending} (7.35,-1.9);
  \draw[call] (7.35,-2.3) -- node[msg, above]{RemoteStart} (9.8,-2.3);
  \draw[return] (9.8,-2.65) -- node[msg, above]{Accepted} (7.35,-2.65);
  \draw[call] (9.8,-3.05) -- node[msg, above]{StartTransaction} (7.35,-3.05);
  \draw[call] (7.35,-3.4) -- node[msg, above]{Événement signé} (2.45,-3.4);
  \draw[return] (2.45,-3.75) -- node[msg, above]{Session créée} (0,-3.75);
  \draw[call] (9.8,-4.2) -- node[msg, above]{MeterValues} (7.35,-4.2);
  \draw[call] (7.35,-4.55) -- node[msg, above]{Mesures} (2.45,-4.55);
  \draw[return] (2.45,-4.9) -- node[msg, above]{Temps réel} (0,-4.9);
  \draw[call] (0,-5.35) -- node[msg, above]{Arrêter} (2.45,-5.35);
  \draw[call] (2.45,-5.75) -- node[msg, above]{RemoteStop} (7.35,-5.75);
  \draw[call] (7.35,-6.15) -- node[msg, above]{Commande} (9.8,-6.15);
  \draw[call] (9.8,-6.55) -- node[msg, above]{StopTransaction} (7.35,-6.55);
  \draw[call] (7.35,-6.9) -- node[msg, above]{Événement signé} (2.45,-6.9);
  \draw[call] (2.45,-7.35) -- node[msg, above]{Capture async} (12.25,-7.35);
  \draw[call] (12.25,-7.8) -- node[msg, above]{Capturer} (4.9,-7.8);
  \draw[return] (4.9,-8.2) -- node[msg, above]{Paiement + reçu} (2.45,-8.2);
\end{tikzpicture}
}
  \caption{Séquence de recharge distante et paiement simulé}
  \label{fig:charging-payment-sequence}
\end{figure}

Une préautorisation réussie ne prouve donc pas que la recharge a commencé. Laravel
crée une commande \texttt{RemoteStartTransaction} chiffrée au repos et limitée dans
le temps. La passerelle la réclame, l'envoie à la connexion active, puis retourne le
résultat. Seul le \texttt{StartTransaction} ultérieur crée la session facturable.
Les \texttt{MeterValues} actualisent l'énergie, la puissance, l'état de charge et
l'estimation de coût. Une limite atteinte produit une commande
\texttt{RemoteStopTransaction}.

L'arrêt manuel du client suit le même mécanisme distant. Le
\texttt{StopTransaction} final fournit le compteur et la raison de fin. Si la borne
se déconnecte avant ce message, la transaction passe en
\texttt{awaiting\_reconciliation} et la session est interrompue sans inventer de
compteur final. Un message tardif peut la réconcilier.

\subsection{Paiement et reçu}

Le paiement applique le patron \textbf{Adapter}. Le domaine dépend de
\texttt{PaymentGateway}, pas d'un fournisseur. Le contrat expose quatre opérations :
préautoriser, capturer, libérer et débiter. Deux implémentations existent :

\begin{itemize}[leftmargin=7mm]
  \item \texttt{SimulatedPaymentAdapter}, rapide et déterministe pour les tests ;
  \item \texttt{WireMockPaymentAdapter}, qui appelle une API HTTP externe simulée.
\end{itemize}

WireMock reproduit quatre résultats : succès, refus métier, délai dépassé et erreur
fournisseur. Les callbacks portent une signature HMAC-SHA256. Laravel vérifie la
signature, enregistre l'identifiant fournisseur une seule fois et protège chaque
opération par une clé d'idempotence. Une répétition du webhook ne peut donc ni
capturer deux fois le paiement ni doubler le chiffre d'affaires.

La préautorisation par défaut est de 30 TND et expire après 15 minutes. Elle est
libérée si la station refuse le démarrage ou si la tentative expire. Après
\texttt{StopTransaction}, le worker capture uniquement le montant final calculé à
partir du tarif figé au démarrage. Le reçu PDF est généré côté serveur à partir de
la session et du paiement persistés. Il peut être prévisualisé ou téléchargé sans
exposer les données d'un autre client.

\subsection{Alertes, interventions et maintenance}

Ces trois notions appartiennent au même cycle opérationnel mais répondent à des
questions différentes :

\begin{description}[leftmargin=31mm,style=nextline]
  \item[Alerte] Quel problème a été détecté, avec quelle gravité et quel SLA ?
  \item[Intervention] Qui traite concrètement ce problème, quand et avec quelles
  preuves ?
  \item[Maintenance] Quelle action préventive ou corrective doit être planifiée,
  éventuellement de façon récurrente ?
\end{description}

\begin{figure}[H]
  \centering
  \resizebox{0.96\textwidth}{!}{\begin{tikzpicture}[
  font=\sffamily\small,
  state/.style={draw=ctline, rounded corners=2mm, thick, fill=white,
    minimum width=28mm, minimum height=12mm, align=center},
  active/.style={state, draw=ctgreen, fill=ctmint},
  flow/.style={-{Latex[length=2.2mm]}, thick, draw=ctgray}
]
  \node[state] (detect) {Événement ou\\détection};
  \node[active, right=7mm of detect] (alert) {Alerte\\ouverte};
  \node[active, right=7mm of alert] (qualify) {Qualification\\et SLA};
  \node[active, right=7mm of qualify] (assign) {Assignation\\technicien};
  \node[active, right=7mm of assign] (work) {Intervention\\en cours};
  \node[state, below=11mm of work] (evidence) {Diagnostic, actions,\\photos et rapport};
  \node[state, left=7mm of evidence] (review) {Contrôle\\opérateur};
  \node[state, left=7mm of review] (close) {Résolution et\\clôture};
  \node[state, below=11mm of qualify] (maintenance) {Plan de maintenance\\ponctuel ou récurrent};

  \draw[flow] (detect) -- (alert);
  \draw[flow] (alert) -- (qualify);
  \draw[flow] (qualify) -- (assign);
  \draw[flow] (assign) -- (work);
  \draw[flow] (work) -- (evidence);
  \draw[flow] (evidence) -- (review);
  \draw[flow] (review) -- (close);
  \draw[flow] (maintenance) -| (assign);
  \draw[flow] (alert) |- node[pos=0.25, left, font=\scriptsize]{préventif} (maintenance);
  \draw[flow] (review) to[bend left=38] (work);
\end{tikzpicture}
}
  \caption{Cycle opérationnel des alertes, interventions et maintenances}
  \label{fig:operations-lifecycle}
\end{figure}

L'opérateur qualifie l'alerte, assigne un technicien et peut créer une intervention.
Le technicien accepte le travail, ajoute notes, diagnostic, photos avant/après et
rapport, puis indique sa résolution. L'opérateur contrôle les preuves et clôture le
suivi. Les changements sont conservés dans des journaux d'événements, et les
notifications SLA sont dédupliquées pour ne pas produire de rappels multiples.

Une maintenance peut exister sans alerte, par exemple un contrôle trimestriel. Le
scheduler génère ses occurrences et peut les transformer en interventions
assignées. À l'inverse, une intervention corrective peut être créée directement
depuis une alerte. Les documents et photos sont conservés sur un disque privé et
servis uniquement après vérification de policy.

\subsection{Tarification, plans et gestion commerciale}

La conception distingue deux marchés qui ne doivent pas être mélangés :

\begin{itemize}[leftmargin=7mm]
  \item le \textbf{tarif de recharge} appartient à l'organisation exploitante et
  calcule le coût d'une session selon le kWh, les frais de session, l'inactivité et
  un minimum ;
  \item le \textbf{plan client} est un abonnement du conducteur à une organisation
  et peut fournir une remise. Un client peut posséder plusieurs abonnements auprès
  d'organisations différentes ;
  \item le \textbf{plan SaaS} est le contrat entre la plateforme et une organisation.
  Il fixe notamment les limites d'employés et de stations.
\end{itemize}

Le tarif effectif est résolu du plus spécifique au plus général : affectation au
connecteur, à la station, puis tarif par défaut de l'organisation. Lors du démarrage,
la session conserve un instantané du tarif et du plan client. Une modification
ultérieure ne change donc pas rétroactivement une recharge en cours.

La demande de démonstration crée un essai commercial. À l'échéance, le cycle de vie
peut entrer en période de grâce, être activé sur une offre ou suspendre les
opérations dépassant les droits. Le super administrateur gère le catalogue SaaS,
les essais et les factures ; l'administrateur d'organisation ne consulte que son
contrat, son usage et ses demandes de changement.

\subsection{Rapports, documents et exports}

Les analyses ne sont pas un écran générique dont seules les valeurs changent. Le
super administrateur suit l'adoption et le portefeuille d'organisations ;
l'administrateur suit les revenus et les capacités ; l'opérateur suit la
disponibilité, les incidents et l'énergie ; le technicien suit ses SLA et résultats
terrain. Les contrôleurs de reporting exposent donc des contrats distincts
\texttt{platform}, \texttt{organization}, \texttt{operations} et \texttt{field}.

Les exports CSV et JSON utilisent les mêmes filtres et le même périmètre que la vue.
Les PDF sont générés côté serveur afin de garantir une identité visuelle stable et
une source vérifiée. Les rapports internes possèdent un émetteur, un destinataire,
un statut, des pièces jointes privées et une version PDF. Ils permettent les
passations entre employés et l'envoi d'un bilan à l'administrateur sans utiliser un
canal externe non audité.
